\section{An arithmetic contrast between two bounded discrete priors}
Continuous smoothing is not interchangeable with adding arbitrary noise
of the same variance. An elementary arithmetic example makes this
particularly explicit. Say $W$ has an \emph{odd-lattice law} if
$W=aM$, where $a>0$, $M$ is symmetric and supported on odd integers,
and $v=\E W^2<\infty$. Its sign pattern
$\eta(W)=\sin(\pi W/(2a))$ takes only the values $\pm1$.

\begin{theorem}[Odd-lattice features retain full parity]
For even $r\ge6$, zero bias, and independent odd-lattice incoming
weights, the full-parity eigenvalue satisfies
\[
                 \lambda_r\ge\frac{4a^4}{\pi^6 r v}.
 \tag{34}
\]
If the greatest common divisor of the supported positive integers $|M|$
is one, then the sharper fixed-law asymptotic is
\[
                 \lambda_r\sim\frac{a^4}{2\pi r v}.
 \tag{35}
\]
Consequently polynomially many such random features support a polynomial
arbitrary-sequence perceptron mistake bound. A finite second moment is
sufficient; Gaussian smoothing and a bounded noise density are not assumed.
\end{theorem}
\begin{proof}
Set $t_0=\pi/(2a)$. The coefficient used for a feature comparator is
$\prod_i\eta(W_i)$, which has square one. Since
\[
 \E[\eta(W)\sin(tW)]=\E\cos((t-t_0)W)=:\psi(t-t_0),
\]
the same cosine-integral and Cauchy--Schwarz argument gives
\[
 \sqrt{\lambda_r}\ge\mu_r:=\frac1\pi\int_0^\infty
                             \psi(t-t_0)^r\frac{dt}{t^2}.
\]
The integrand is nonnegative. On
$[t_0-(rv)^{-1/2},t_0]$, one has
$\psi(t-t_0)\ge1-v(t-t_0)^2/2\ge1-1/(2r)$.
This interval is positive because $v\ge a^2$ and $r\ge6$.
Using $\psi^r\ge1/2$ and $t^{-2}\ge t_0^{-2}$ yields
$\mu_r\ge2a^2/(\pi^3\sqrt{rv})$, proving (34).

For (35), Minkowski supplies the upper bound
\[
 \sqrt{\lambda_r}\le\frac1\pi\int_0^\infty
                \{\E\sin^2(tW)\}^{r/2}\frac{dt}{t^2}.
\]
The upper integrand reaches its maximum exactly at
$t_m=(2m+1)\pi/(2a)$. To verify this assertion, equality requires
$t aM$ to be an odd multiple of $\pi/2$ for every supported integer.
The gcd assumption and an integer combination of a finite subset of the
support force $2at/\pi$ to be an integer; its product with every odd
$M$ must be odd, so that integer is odd.
Both integrands have these peaks and are periodic, after their even
powers are taken, with period $\pi/a$. At every peak, finite variance gives
\[
 \psi(t-t_0)^r=e^{-rv(t-t_m)^2/2}(1+o(1)),\quad
 \{\E\sin^2(tW)\}^{r/2}=e^{-rv(t-t_m)^2/2}(1+o(1))
\]
on the $r^{-1/2}$ scale. Compactness of a period gives a fixed gap
away from the peaks; the local quadratic expansion supplies a Gaussian
envelope near each peak. Near the original endpoint zero, use
$\E\sin^2(tW)\le vt^2$ to control the apparent $t^{-2}$ singularity.
The remaining periods are summable because their weights are $O(m^{-2})$.
Dominated period-by-period Laplace integration therefore makes both root
bounds asymptotic to
\[
 \frac1\pi\sqrt{\frac{2\pi}{rv}}
          \sum_{m\ge0}t_m^{-2}
 =\frac{a^2}{\sqrt{2\pi rv}},
\]
using $\sum_{m\ge0}(2m+1)^{-2}=\pi^2/8$. Squeezing and squaring
prove (35).

For finite features the comparator uses the original incoming weights
and the bounded coefficient $\prod_i\eta(W_i)$. Symmetry of $W$ and
oddness of $\eta$ give its expected output $\mu_r\chi(x)$ at every
input. The preceding expectation-symmetrization argument applies with
$1+\sigma^2$ replaced by $v$. Equation (34) makes all its sufficient
width, norm, radius, and perceptron bounds polynomial in $r$ at a fixed law.
\end{proof}

For example compare these two mean-zero, variance-one, bounded four-point
initializations:
\[
 W_{\rm odd}\sim\operatorname{Unif}\{-3,-1,1,3\}/\sqrt5,
 \qquad
 W_{\rm mixed}\sim\operatorname{Unif}\{-2,-1,1,2\}/\sqrt{5/2}.
\]
For the first, (35) gives $\lambda_r\sim1/(50\pi r)$.
For the second, scale does not affect the sine gap, and
\[
 \sup_t\frac{\sin^2t+\sin^2(2t)}2=\frac{25}{32}.
\]
Indeed the expression is $3/4-c/4-c^2/2$ with $c=\cos2t$, whose
maximum is attained at $c=-1/4$. Equation (2) therefore gives
\[
 \lambda_r(W_{\rm mixed})\le\frac1{\pi^2}
       \left(\frac r{r-1}\right)^2(25/32)^{r-1}.       \tag{36}
\]
The exponential base is sharp for this second example:
\[
 \lambda_r(W_{\rm mixed})\sim
       \frac{22}{25\pi\sqrt3\,r}(25/32)^r.             \tag{37}
\]
Here is a direct two-dimensional Laplace proof, including the constant.
Work first with the unnormalized law on $\{\pm1,\pm2\}$. Fubini in
the absolutely convergent Fourier integral gives
\[
 \lambda_r=\frac1{\pi^2}\int_0^\infty\int_0^\infty
 \left\{\frac{\sin t\sin u+\sin(2t)\sin(2u)}2\right\}^{r}
                                      \frac{dt\,du}{t^2u^2}.
\]
Let $t_*\in(0,\pi/2)$ satisfy $\cos(2t_*)=-1/4$, and set
$\kappa=25/32$. Cauchy--Schwarz and the preceding maximum calculation
show that the absolute covariance reaches $\kappa$ precisely when
$t,u$ belong to the same one of these two families, modulo $2\pi$:
\[
 \mathcal A=\{\pm t_*\},\qquad
 \mathcal B=\{\pi\pm t_*\}.
\]
Only positive representatives are used in the integral. The Hessian of
minus the logarithm of the absolute covariance at each such pair has
diagonal entries $14/5$ and off-diagonal entries $\pm2/5$; its determinant
is $192/25$. For example, at $(t_*,t_*)$ the underlying second moments
are $\E[W^2\sin^2(t_*W)]=35/16$ and
$\E[W^2\cos^2(t_*W)]=5/16$.
Thus each pair contributes the common Laplace factor
$5\kappa^r/(4\pi\sqrt3\,r)$ times $t^{-2}u^{-2}$.
Periodicity and a fixed gap off these finitely many peaks in a period
justify the two-dimensional localization. The period weights are
summable; near either zero axis the elementary bound
$|\sin(tW)|\le|tW|$ controls its apparent singularity.

The two positive resonance sums are
\[
 S_A=\sum_{n\in\mathbb Z}(t_*+2\pi n)^{-2}
       =\frac1{4\sin^2(t_*/2)},\qquad
 S_B=\frac1{4\cos^2(t_*/2)}.
\]
Since $\cos^2t_*=3/8$, $S_A^2+S_B^2=44/25$.
The unnormalized kernel is therefore asymptotic to
$11\kappa^r/(5\pi\sqrt3\,r)$. Dividing by its variance $5/2$
proves (37). An accompanying exact-integer calculation averages over the
binomial count of magnitude-two weights and checks this asymptotic;
that calculation is a diagnostic, not the Laplace proof.

Thus bounded discrete priors with the same first two moments can have
inverse-polynomial versus exponentially small parity features.
The distinction here is the arithmetic support and its common sine
resonances. This is an initial frozen-feature result; it does not classify
all discrete priors or prove a corresponding full-network optimization
separation.

\paragraph{Independent hidden biases need not destroy the arithmetic advantage.}
The zero-bias condition can be relaxed for the odd-lattice example. Let $B$
be a fixed independent bias law with $\E B^2<\infty$. Define the bounded triangular wave
\[
 T_a(b)=\frac8{\pi^2}\sum_{m\ge0}
       \frac{\cos((2m+1)\pi b/(2a))}{(2m+1)^2}.
\]
It is $4a$-periodic and equals $1-|b|/a$ on $|b|\le2a$.
Under the gcd assumption of (35), the biased kernel has
\[
 \lambda_r(B)\sim\frac{a^4}{2\pi r v}\,\E T_a(B)^2
 \quad\text{whenever }\E T_a(B)^2>0.                 \tag{38}
\]
Equivalently, $r\lambda_r(B)$ always converges to the displayed constant
without the factor $1/r$, including when that constant is zero.
Every nondegenerate continuous bias law satisfies the positivity condition:
the zeros of $T_a$ are only the odd multiples of $a$.
In particular, the inverse-polynomial versus exponential comparison of the
two four-point incoming laws remains valid with the same fixed Gaussian
hidden-bias distribution. The mixed-law bound (36) holds for any independent
bias by the original Fourier upper bound.

To verify (38), first fix a real bias $b$. The squared Fourier integral
acquires the factor $\cos(tb)\cos(ub)$. For the odd-lattice law, the
maximal absolute covariance is one, attained at every pair
$(t_m,t_n)$ of the odd half-periods. The Hessian of minus its logarithm
is $\operatorname{diag}(v,v)$: each single-weight cosine vanishes at a
resonance, so the cross derivative is zero. Two-dimensional Laplace
integration, with the same summable period weights, therefore gives
\[
 r\lambda_r(b)\longrightarrow
 \frac2{\pi v}\left\{\sum_{m\ge0}
                          \frac{\cos(t_mb)}{t_m^2}\right\}^2
              =\frac{a^4}{2\pi v}T_a(b)^2.
\]
The bound $\lambda_r(b)\le J_r^2$ from the odd-lattice proof is
independent of $b$, and $rJ_r^2$ is bounded. Dominated convergence
justifies averaging over the stated fixed law of $B$.
The triangular-wave formula follows from the absolutely convergent
odd-cosine series, or by integrating the square-wave sine series.

Actual finite features still have a polynomial perceptron guarantee. Use the bounded comparator coefficient
$T_a(B)\prod_i\eta(W_i)$. The same one-dimensional peak calculation
gives its mean full-parity coefficient asymptotic in absolute value to
$a^2\E T_a(B)^2/\sqrt{2\pi rv}$; the coefficient is thus nonzero and
inverse polynomial at all sufficiently large even dimensions.
Symmetrization uses the augmented incoming second moment $rv+\E B^2$,
and the previous feature-radius argument applies. This conclusion assumes
the stated fixed bias law, and does not infer a positive leading constant
for a bias supported entirely on the triangular wave's zeros.

\begin{theorem}[A cancelled leading term can still be polynomial]
Assume the odd-lattice law has finite support and the gcd in (35) is one.
At the fixed hidden bias $b=a$, the full-parity kernel satisfies
\[
             \lambda_r(a)\sim\frac{a^8}{2\pi r^3v^3}.
 \tag{39}
\]
Polynomially many actual features again support a polynomial
arbitrary-sequence perceptron bound. In particular, for the variance-one
law on $\{\pm1,\pm3\}/\sqrt5$ and bias $1/\sqrt5$, the eigenvalue
is asymptotic to $1/(1250\pi r^3)$. Thus the zero of the leading
triangular-wave term in (38) need not imply exponential suppression.
\end{theorem}
\begin{proof}
Let $\epsilon_r\in\{\pm1\}$ be the deterministic phase in the
full-parity cosine integral, and write
$\eta_* =\prod_i\eta(w_i)$ and $V_r=\sum_iw_i^2$.
For each fixed vector of odd-lattice weights, periodize the full real-line
integral around $t_m=(m+1/2)\pi/a$, $m\in\mathbb Z$. On
$|u|<\pi/(2a)$ the sine product is
$\eta_*\prod_i\cos(w_i u)$ at every such period, because $r$ is even.
The elementary cosecant partial-fraction identity, differentiated once,
gives
\[
 \sum_{m\in\mathbb Z}\frac{\cos(a(t_m+u))}{(t_m+u)^2}
 =a^2\tan^2(au).
\]
For completeness, use $\cos(a(t_m+u))=-(-1)^m\sin(au)$ together with
\[
 \sum_m(-1)^m(t_m+u)^{-2}=-a^2\sec(au)\tan(au).
\]
Absolute convergence away from the period endpoints and the order-$r$
zeros of the sine product at those endpoints justify periodization.
Consequently the actual finite-cube ReLU coefficient has the exact formula
\[
 G_{[r]}(w,a)=\frac{\epsilon_r\eta_*a^2}{2\pi}
   \int_{-\pi/(2a)}^{\pi/(2a)}
          \tan^2(au)\prod_{i=1}^r\cos(w_i u)\,du.
 \tag{40}
\]

With probability $1-e^{-cr}$, every supported magnitude occurs with
at least half its expected frequency. On this event the gcd assumption
provides a uniform exponential gap for the absolute cosine product away
from $u=0$. The apparent endpoint poles are harmless: retain two cosine
factors to cancel $\tan^2(au)$, and bound the remaining factors using
those frequencies. Near zero, bounded support gives, uniformly in these
weight vectors,
\[
 \prod_i\cos(w_i u)=\exp\{-V_r u^2/2+O(ru^4)\},
 \qquad a^2\tan^2(au)=a^4u^2+O(u^4).
\]
A Gaussian envelope and $V_r\asymp r$ permit dominated Laplace
integration. Formula (40) becomes
\[
 G_{[r]}(w,a)=\epsilon_r\eta_*\,
       \frac{a^4}{\sqrt{2\pi}\,V_r^{3/2}}(1+o(1)),
\]
uniformly on the stated frequency event. Moreover $V_r/r\to v$ in
probability, and $V_r\ge a^2r$. The coefficient is $O(\sqrt r)$
on its exponentially unlikely complement, by Cauchy--Schwarz in the
original cube average. Squaring and averaging prove (39).

The same calculation without squaring shows that the bounded comparator
coefficient $\epsilon_r\prod_i\eta(W_i)$ has expected output
$\mu_r\chi(x)$ at every cube input, with
$\mu_r\sim a^4/(\sqrt{2\pi}(rv)^{3/2})$. Sign symmetry gives this
identity simultaneously for all inputs, even though the hidden bias is
fixed at $a$. The finite-feature symmetrization, radius, and perceptron
argument therefore applies again. For a fixed law and fixed confidence,
it gives sufficient width $O(r^5)$ and mistake count $O(r^5)$.
\end{proof}
